% HLE and AIME macro analysis sections for AlphaDiana paper
% Requires: \usepackage{tcolorbox} in preamble
% \findingbox command and \newcounter{finding} are defined in gpqa_macro_analysis.tex

\subsection{HLE}
\label{appendix:macro-hle}

We evaluate three systems on the HLE \texttt{multipleChoice} subset using
Qwen3.5-27B: DirectLLM ($n=323$), OpenCode ($n=591$), and ZeroClaw ($n=591$).
OpenClaw is not evaluated on this benchmark.
Table~\ref{tab:hle-overview} reports accuracy, output length, entropy, and
tool-use rates per system.
Terms carry the same definitions as in Section~\ref{appendix:macro-gpqa}.

The benchmark inverts the GPQA-Diamond harness ranking.
DirectLLM leads at 23.5\% versus 13.9\% and 14.9\% for the two agent
harnesses, a sign reversal absent in GPQA.
Three contrasts define the difference.
First, DirectLLM wrong trajectories hit the 32k-token context limit at 22\%,
meaning the model generates extended reasoning chains and still fails; the
agent harnesses have compressed budgets and shorter trajectories at every
outcome.
Second, the tool-use rate on OpenCode is 23\% on correct and 24\% on wrong:
the 3.6$\times$ outcome-conditioned gap from GPQA has collapsed to near zero.
Third, the W/C token ratio for ZeroClaw is 0.9$\times$, meaning correct
trajectories are \emph{longer} than wrong ones; the token-inflation signature
from GPQA Finding~1 does not apply here.

Output entropy also changes character on HLE.
DirectLLM entropy is flat across outcomes (0.34 correct vs.\ 0.33 wrong),
carrying no diagnostic content.
OpenCode correct trajectories have \emph{higher} mean entropy (0.42 vs.\
0.33), the inverse of the low-entropy collapse observed for OpenClaw on GPQA;
correct HLE trajectories are more exploratory, not more repetitive.
ZeroClaw shows a slight downward shift on wrong (0.31 vs.\ 0.34), with no
length signal.
Entropy is therefore a benchmark-conditional signal: the patterns identified
on GPQA do not transfer to HLE.

\begin{table}[t]
\centering
\small
\caption{%
  System-level overview on HLE \texttt{multipleChoice} ($n_{\text{shared}}=323$,
  Qwen3.5-27B; OpenCode and ZeroClaw evaluated on $n=591$ including tasks
  outside the DirectLLM set).
  Tok\,(C)/Tok\,(W) = mean output tokens on correct/wrong trajectories.
  W/C = Tok\,(W)/Tok\,(C).
  Ent\,(C/W) = mean per-token log-probability entropy.
  Tool\,\% = tool use rate, correct/wrong.
  Loss\,\% = harness-specific failure rate over the 76 DirectLLM-correct tasks.
}
\label{tab:hle-overview}
\setlength{\tabcolsep}{3pt}
\begin{tabular}{lrrrrrrr}
\toprule
System & Acc\,\% & Tok (C) & Tok (W) & W/C & Ent (C\,/\,W) & Tool\,\% (C\,/\,W) & Loss\,\% \\
\midrule
DirectLLM &  23.5 & 12{,}175 & 14{,}710 & $1.2{\times}$ & 0.34\,/\,0.33 & 0\,/\,0    & N/A  \\
OpenCode  &  13.9 &  1{,}466 &  7{,}373 & $5.0{\times}$ & 0.42\,/\,0.33 & 23\,/\,24  & 67.1 \\
ZeroClaw  &  14.9 &  2{,}340 &  2{,}073 & $0.9{\times}$ & 0.34\,/\,0.31 & 0\,/\,0    & 65.8 \\
\midrule
Oracle    &  \textbf{33.1} & \multicolumn{6}{c}{(per-task best available system)} \\
\bottomrule
\end{tabular}
\end{table}

\paragraph{Action Distribution.}
Table~\ref{tab:hle-action-rates} reports outcome-conditioned action rates.
DirectLLM rates reflect raw-text keyword detection and are heuristic;
the only precise column is tool use (0\% for DirectLLM, confirmed from
event logs).
OpenCode tool use is flat across outcomes (23\% correct, 24\% wrong): the
outcome-predictive tool activation signal from GPQA-Diamond does not appear
on HLE.
ZeroClaw wrong trajectories over-emit recovery actions relative to correct
ones (36\% vs.\ 30\%), the only outcome-conditioned action gap across any
(harness, benchmark) combination on HLE.

\begin{table}[t]
\centering
\small
\caption{%
  Outcome-conditioned action rates on HLE ($n_{\text{shared}}=323$ for
  DirectLLM; $n=590/591$ for OpenCode/ZeroClaw).
  Each value is the percentage of trajectories containing at least one event
  of that action type.
  \texttt{verify}, \texttt{plan}, \texttt{recover}, \texttt{reason} for
  DirectLLM are keyword-heuristic counts over raw output; \texttt{tool\_use}
  is from event logs (precise for all harnesses).
}
\label{tab:hle-action-rates}
\setlength{\tabcolsep}{4pt}
\begin{tabular}{llrcccccc}
\toprule
Harness & Outcome & $n$ & Ans.\,\% & Reason\,\% & Tool\,\% & Verify\,\% & Plan\,\% & Recover\,\% \\
\midrule
DirectLLM & correct &  76 & 100 & 78 &  0 &  4 &  0 &  7 \\
DirectLLM & wrong   & 247 & 100 & 65 &  0 &  5 &  0 &  7 \\
\midrule
OpenCode  & correct &  82 & 100 & 59 & 23 & 26 &  7 &  9 \\
OpenCode  & wrong   & 508 & 100 & 50 & 24 & 28 & 11 &  9 \\
\midrule
ZeroClaw  & correct &  88 & 100 & 50 &  0 & 19 & 17 & 30 \\
ZeroClaw  & wrong   & 503 & 100 & 43 &  0 & 19 & 16 & \textbf{36} \\
\bottomrule
\end{tabular}
\end{table}

\findingbox{On knowledge-intensive benchmarks with no executable solution
path, extended per-sample reasoning budget is the primary accuracy driver;
harnesses that compress the per-trajectory budget replace thinking time with
tool overhead and lose accuracy relative to unconstrained direct inference.}

HLE questions require broad specialized knowledge across dozens of academic
subfields; unlike GPQA-Diamond's numerical subproblems, they offer no
tractable executable path.
DirectLLM, unconstrained, allocates 12k tokens to correct trajectories and
14.7k to wrong ones; 22\% of wrong trajectories exhaust the 32k-token limit
entirely.
Agent harnesses operate within per-step budgets that produce trajectories
roughly one-sixth as long on correct instances, compressing the effective
reasoning window without providing a compensating tool benefit.
Table~\ref{tab:hle-goal-loss} shows that ZeroClaw's W/C ratio inverts to
0.9$\times$ because both correct and wrong trajectories are forced into the
same short budget; length no longer discriminates outcomes.

\textbf{Case: hle\_155 (linear algebra, rank-1 approximation).}
DirectLLM works through a 21k-token derivation and answers correctly.
OpenCode exhausts its per-step budget and returns no answer.
ZeroClaw commits to the wrong option in 1,853 tokens, less than one-tenth of
DirectLLM's allocation.

\begin{table}[t]
\centering
\small
\caption{%
  Trajectory-level signals by system and outcome on HLE.
  Mean tokens = mean output length.
  Cap\,\% = fraction of trajectories reaching the 32k-token context limit.
  DirectLLM cap events reflect extended reasoning that exceeds the window, not
  deliberation looping.
}
\label{tab:hle-goal-loss}
\setlength{\tabcolsep}{4pt}
\begin{tabular}{llrrr}
\toprule
System & Outcome & $n$ & Mean tokens & Cap\,\% \\
\midrule
DirectLLM & correct &  76 & 12{,}175 &  5 \\
DirectLLM & wrong   & 247 & 14{,}710 & 22 \\
\midrule
OpenCode  & correct &  82 &  1{,}466 &  0 \\
OpenCode  & wrong   & 508 &  7{,}373 &  9 \\
\midrule
ZeroClaw  & correct &  88 &  2{,}340 &  0 \\
ZeroClaw  & wrong   & 503 &  2{,}073 &  0 \\
\bottomrule
\end{tabular}
\end{table}

\findingbox{When task structure provides no executable solution path, tool
activation rate becomes uniform across outcomes; the presence of tool
invocations neither predicts nor explains correct answers.}

OpenCode invokes tools in 23\% of correct and 24\% of wrong trajectories,
a near-zero outcome-conditioned gap.
Unlike GPQA-Diamond, where tool use splits sharply by outcome because some
subdomains support numerical computation, HLE questions require expert
knowledge in philosophy, linguistics, theoretical mathematics, and
experimental science; no code execution or web search can substitute for
that knowledge.
Tool calls on HLE are exploratory rather than anchoring: the model searches
for information that would confirm a chosen answer rather than computing the
answer from first principles.
The tool-use gap is therefore a property of the benchmark's computational
structure, not of the harness design.

\textbf{Case: hle\_68 (Japanese pitch accent).}
OpenCode invokes a search tool and commits to the wrong accent class in 551
tokens.
DirectLLM enumerates the four pitch accent categories systematically across
32k tokens and arrives at the correct class through exhaustive elimination.
The task requires recalling and applying a specific phonological rule; no
tool result substitutes for that knowledge.

\findingbox{On benchmarks where base-model accuracy is low, the paired-outcome
decomposition reveals an asymmetry between failure types: harness-induced
losses are predominantly premature commitments reached with a fraction of the
base model's reasoning budget, not format extraction failures.}

The \emph{harness loses} row of Table~\ref{tab:hle-failure-sig} isolates
harness-induced damage.
\emph{Malformed} means the extracted answer does not match any valid option
letter.
\emph{TokRatio} is the harness token count divided by the DirectLLM token
count on the same task; values below 1 mean the harness used less computation.
Unlike GPQA-Diamond, where OpenClaw regressions featured 89\% malformed and
83\% no-boxed answers from context overflow, HLE regressions are mostly
extractable: OpenCode's 36\% malformed and ZeroClaw's 24\% are the primary
format issues, while TokRatio values well below 1 confirm that the harness
answered with far less deliberation than DirectLLM needed.
Both harnesses commit quickly; neither spirals.

\textbf{Case: hle\_583 (causal regression with GMM).}
DirectLLM identifies the correct variable after 24k tokens of systematic
elimination.
ZeroClaw outputs \texttt{3} (a malformed answer) in 590 tokens.
OpenCode selects the wrong letter in 3,459 tokens.
Both harnesses commit at a fraction of DirectLLM's budget and arrive at
wrong or invalid answers; the failure is premature commitment, not format
extraction.

\begin{table}[t]
\centering
\small
\caption{%
  Paired-outcome failure signatures on HLE.
  Paired on the 323 DirectLLM-evaluated tasks.
  Malformed\,\% = extracted answer does not match a valid option letter.
  TokRatio = harness token count / DirectLLM token count on the same task.
  HLE uses letter-based answers; no-boxed rate is 0\% for all cells and
  is omitted.
}
\label{tab:hle-failure-sig}
\setlength{\tabcolsep}{3.5pt}
\begin{tabular}{llrrr}
\toprule
Harness & Outcome & $n$ & Malformed\,\% & TokRatio \\
\midrule
\multirow{4}{*}{OpenCode}
 & both correct  &  25 &  0.0 & 0.26 \\
 & harness gains &  18 &  0.0 & 0.12 \\
 & harness loses &  51 & 36.0 & 0.86 \\
 & both wrong    & 229 & 29.7 & 0.46 \\
\midrule
\multirow{4}{*}{ZeroClaw}
 & both correct  &  26 &  0.0 & 0.34 \\
 & harness gains &  25 &  0.0 & 0.22 \\
 & harness loses &  50 & 24.0 & 0.21 \\
 & both wrong    & 222 & 31.5 & 0.22 \\
\bottomrule
\end{tabular}
\end{table}

\findingbox{When base-model accuracy is low, harness overhead concentrates on
already-correct instances; both recovery rate and oracle ceiling are bounded
by the fraction of tasks any system can solve, not by harness design quality.}

The \emph{recovery rate} and \emph{loss rate} definitions follow
Section~\ref{appendix:macro-gpqa}.
Table~\ref{tab:hle-oracle} shows that both harnesses achieve strongly
negative net gains: OpenCode loses 67\% of DirectLLM-correct tasks while
recovering only 7\% of DirectLLM-failed tasks; ZeroClaw loses 66\% and
recovers 10\%.
The loss-to-recovery ratio is roughly 9:1 on HLE, compared to 1:2 on
GPQA-Diamond (where ZeroClaw recovered 33\% while losing 9\%).
The oracle ceiling of 33.1\% is only 9.6 percentage points above DirectLLM,
compared to an 8.1-point gap on GPQA with a much higher absolute ceiling.
On HLE, the ceiling is not constrained by harness routing but by the fraction
of problems any of the three systems can solve at all.

\begin{table}[t]
\centering
\small
\caption{%
  Oracle harness router ceiling on HLE (paired on $n=323$ DirectLLM tasks).
  Recovery\,\% = harness-only success rate over the 247 DirectLLM-wrong tasks.
  Loss\,\% = harness-specific failure rate over the 76 DirectLLM-correct tasks.
  Oracle selects, per task, any harness that answers correctly when one exists.
}
\label{tab:hle-oracle}
\setlength{\tabcolsep}{4pt}
\begin{tabular}{lrrr}
\toprule
System & Accuracy\,\% & Recovery\,\% & Loss\,\% \\
\midrule
DirectLLM & 23.5 & N/A  & N/A \\
OpenCode  & 13.9 &  7.3 & 67.1 \\
ZeroClaw  & 14.9 & 10.1 & 65.8 \\
\midrule
Oracle    & \textbf{33.1} & 12.6 & N/A \\
\bottomrule
\end{tabular}
\end{table}


% ---------------------------------------------------------------------------
\paragraph{ARM Cognitive-Mode Analysis on HLE.}
% ---------------------------------------------------------------------------

The ARM (Agentic Reasoning Mode) taxonomy defined in
Section~\ref{appendix:macro-gpqa} classifies every assistant-generated segment
into one of seven cognitive modes: principle-based deduction (PD), systematic
elimination (SE), intuitive commitment (IC), uncertainty navigation (UN),
recovery-replanning (RR), and undifferentiated reasoning (reason).
We apply the same keyword classifiers to HLE trajectories.
Table~\ref{tab:hle-arm-modes} reports the outcome-conditioned ARM distribution
for DirectLLM and OpenCode on HLE, alongside the GPQA-Diamond values for
contrast.

\begin{table}[t]
\centering
\small
\caption{%
  ARM cognitive-mode distribution by benchmark and outcome.
  HLE values from the same keyword classifiers applied to HLE trajectory text.
  $\Delta$PD = PD(wrong) $-$ PD(correct).
  GPQA values reproduced from Section~\ref{appendix:macro-gpqa} for contrast.
}
\label{tab:hle-arm-modes}
\setlength{\tabcolsep}{2.5pt}
\begin{tabular}{llrrrrrrr}
\toprule
Bench. & Harness & Out. & PD & reason & IC & UN & RR & $\Delta$PD \\
\midrule
\multirow{6}{*}{GPQA}
& \multirow{2}{*}{OpenClaw}
  & C & 0.538 & 0.372 & 0.049 & 0.016 & 0.017 & \multirow{2}{*}{+0.070} \\
& & W & 0.608 & 0.291 & 0.026 & 0.043 & 0.010 & \\
\cmidrule{2-8}
& \multirow{2}{*}{OpenCode}
  & C & 0.353 & 0.587 & 0.026 & 0.009 & 0.019 & \multirow{2}{*}{+0.145} \\
& & W & 0.498 & 0.449 & 0.000 & 0.042 & 0.011 & \\
\cmidrule{2-8}
& \multirow{2}{*}{ZeroClaw}
  & C & 0.280 & 0.664 & 0.023 & 0.011 & 0.021 & \multirow{2}{*}{+0.126} \\
& & W & 0.406 & 0.504 & 0.011 & 0.032 & 0.044 & \\
\midrule
\multirow{4}{*}{HLE}
& \multirow{2}{*}{DirectLLM}
  & C & 0.125 & 0.417 & 0.082 & 0.099 & 0.272 & \multirow{2}{*}{\textbf{$-$0.019}} \\
& & W & 0.106 & 0.423 & 0.074 & 0.135 & 0.259 & \\
\cmidrule{2-8}
& \multirow{2}{*}{OpenCode}
  & C & 0.162 & 0.142 & 0.132 & 0.201 & 0.364 & \multirow{2}{*}{\textbf{$-$0.017}} \\
& & W & 0.144 & 0.153 & 0.152 & 0.238 & 0.310 & \\
\bottomrule
\end{tabular}
\end{table}

\findingbox{The PD-sink over-reasoning signature from GPQA-Diamond does not
appear on HLE.
ARM failure modes are determined by the interaction between benchmark domain
structure and harness architecture, not by a fixed model-level cognitive bias.}

On GPQA-Diamond, PD is the dominant mode (28\% to 61\% of segments) and PD
over-expression is the strongest single-mode predictor of failure, with
wrong$-$correct PD gaps of +7.0 to +14.5~pp across harnesses.
On HLE, PD is a minority mode (11\% to 16\%), and the wrong$-$correct PD gap
is \emph{negative} for both DirectLLM ($-$1.9~pp) and OpenCode ($-$1.7~pp):
wrong HLE trajectories allocate \emph{less} of their cognitive budget to
deduction, not more.

Three structural differences drive the divergence.
First, GPQA-Diamond questions are hard-science problems where PD is the natural
reasoning strategy; HLE questions span philosophy, linguistics, theoretical
mathematics, and experimental science, where no single derivational framework
applies.
Second, the dominant non-reason mode on HLE is recovery: RR rates of 0.26 to
0.36, roughly 10$\times$ the GPQA baseline.
The model is constantly reconsidering its approach under genuine domain
uncertainty, not spiraling in a single mode.
Third, arm entropy on HLE DirectLLM is substantially higher (1.14 to 1.23) than
on any GPQA harness (0.40 to 0.81), confirming that mode diversity reflects
task diversity rather than failure.

The implication is that a harness designed to prevent the PD spiral on GPQA
(e.g., ZeroClaw's fixed template) may have no effect on HLE, where the failure
mode is premature commitment under knowledge uncertainty rather than
over-deduction.

% ---------------------------------------------------------------------------
\paragraph{Context rotting on HLE is outcome-invariant.}
% ---------------------------------------------------------------------------

Table~\ref{tab:hle-context-rotting} extends the context-rotting diagnostic from
GPQA (Section~\ref{appendix:macro-gpqa}) to HLE.
Head entropy is the mean per-token log-probability entropy over the first
quartile of tokens; tail entropy is the same over the last quartile.
$\Delta$H = tail $-$ head.

\begin{table}[t]
\centering
\small
\caption{%
  Context rotting signals on HLE, by harness and outcome.
  Head = mean entropy over first quartile of tokens.
  Tail = mean over last quartile.
  $\Delta$H = tail $-$ head.
  Low-ent\,\% = fraction of tokens in the bottom quartile of the per-trajectory
  entropy distribution.
}
\label{tab:hle-context-rotting}
\setlength{\tabcolsep}{3pt}
\begin{tabular}{llrrrrr}
\toprule
Harness & Outcome & Tokens & Head & Tail & $\Delta$H & Low-ent\,\% \\
\midrule
\multirow{2}{*}{DirectLLM}
  & C & 12{,}175 & 0.402 & 0.271 & $-$0.131 & 0.249 \\
  & W & 14{,}710 & 0.407 & 0.260 & $-$0.147 & 0.244 \\
\midrule
\multirow{2}{*}{OpenCode}
  & C &  1{,}466 & 0.372 & 0.328 & $-$0.044 & 0.250 \\
  & W &  7{,}373 & 0.324 & 0.262 & $-$0.061 & 0.250 \\
\midrule
\multirow{2}{*}{ZeroClaw}
  & C &  2{,}340 & 0.334 & 0.272 & $-$0.062 & 0.250 \\
  & W &  2{,}287 & 0.338 & 0.271 & $-$0.067 & 0.250 \\
\bottomrule
\end{tabular}
\end{table}

\findingbox{Entropy decline on HLE is outcome-invariant: both correct and wrong
trajectories lose approximately 0.14 nats from head to tail.
Unlike GPQA OpenClaw, whose capped trajectories collapse to tail entropy 0.001,
HLE DirectLLM trajectories end at 0.26, well above the threshold where answer
extraction becomes unreliable.}

Table~\ref{tab:hle-context-rotting} reveals a pattern distinct from the three
GPQA rotting regimes.
HLE DirectLLM trajectories start at head entropy approximately 0.40 and end at
tail entropy approximately 0.26 regardless of correctness.
This is a natural consequence of long-form reasoning: as the model narrows
toward a conclusion, its token distribution sharpens.
The decline is not diagnostic of failure.

Critically, HLE DirectLLM tail entropy (0.26) remains well above the
catastrophic threshold observed on GPQA OpenClaw capped trajectories (0.001).
At tail entropy 0.26, the model is still deliberating; at 0.001, it is cycling
through deterministic repetition.
The 22\% of HLE DirectLLM wrong trajectories that hit the 32K context limit
(Table~\ref{tab:hle-goal-loss}) do so through extended reasoning that
progressively sharpens toward a conclusion, not through the PD sink that
destroys OpenClaw trajectories.

HLE OpenCode and ZeroClaw show minimal $\Delta$H ($-$0.04 to $-$0.07),
consistent with their compressed token budgets preventing any significant
entropy evolution from head to tail.

% ---------------------------------------------------------------------------
\paragraph{Recovery on HLE is productive exploration, not mode oscillation.}
% ---------------------------------------------------------------------------

Table~\ref{tab:hle-rr-context} places HLE RR rates alongside arm entropy and
PD rate for interpretability.

\begin{table}[t]
\centering
\small
\caption{%
  Recovery (RR) mode rates in context on HLE versus GPQA.
  Arm entropy = cognitive mode diversity.
  GPQA values use the four-way paired decomposition; HLE values use simple
  correct/wrong.
}
\label{tab:hle-rr-context}
\setlength{\tabcolsep}{3pt}
\begin{tabular}{llrrrrr}
\toprule
Bench. & Harness & Outcome & RR rate & Arm ent. & PD rate & Tokens \\
\midrule
\multirow{6}{*}{GPQA}
& \multirow{2}{*}{OpenClaw}
  & regression & 0.005 & 0.261 & 0.643 & 30{,}451 \\
& & both correct & 0.018 & 0.476 & 0.519 &  1{,}932 \\
\midrule
& \multirow{2}{*}{ZeroClaw}
  & regression & 0.044 & 0.825 & 0.427 &  1{,}913 \\
& & both correct & 0.020 & 0.650 & 0.271 &  1{,}982 \\
\midrule
\multirow{4}{*}{HLE}
& \multirow{2}{*}{DirectLLM}
  & C & 0.272 & 1.230 & 0.125 & 12{,}175 \\
& & W & 0.259 & 1.136 & 0.106 & 14{,}710 \\
\cmidrule{2-7}
& \multirow{2}{*}{OpenCode}
  & C & 0.364 & 0.166 & 0.162 &  1{,}466 \\
& & W & 0.310 & 0.128 & 0.144 &  7{,}373 \\
\bottomrule
\end{tabular}
\end{table}

\findingbox{The same RR rate carries opposite diagnostic meaning on different
benchmarks.
On HLE, RR rates of 0.26 to 0.36 with arm entropy 1.1 to 1.2 and PD below 0.16
describe productive exploration under genuine domain uncertainty.
On GPQA ZeroClaw, RR rates of 0.04 with arm entropy above 0.8 and PD above 0.40
describe mode oscillation without convergence.}

On HLE, RR is the dominant non-reason cognitive mode, exceeding PD by roughly
2$\times$.
Arm entropy is high (1.1 to 1.2 for DirectLLM), and PD rates are at their
lowest (0.11 to 0.16).
This profile describes a model that is genuinely navigating uncertainty across
diverse academic subfields: it frequently reconsiders its approach, draws on
multiple cognitive strategies, and does not default to deduction under ambiguity.

On GPQA ZeroClaw, RR rates are only 0.02 to 0.04, but when they appear on wrong
trajectories, they are embedded in high arm entropy (0.76 to 0.83) and elevated
PD (0.39 to 0.43).
The model oscillates between modes rather than executing a focused correction.
On GPQA OpenClaw regression, RR is near zero (0.005): the model cannot recover
at all because it cannot exit the PD sink.

The triple (RR rate, arm entropy, PD rate) jointly determines whether recovery
is productive: high RR with high arm entropy and low PD signals exploration;
high RR with high arm entropy and high PD signals oscillation; low RR with low
arm entropy and high PD signals collapse.

% ---------------------------------------------------------------------------
\paragraph{Cross-benchmark action rates: the tool-use gap is
benchmark-conditional.}
% ---------------------------------------------------------------------------

Table~\ref{tab:hle-cross-actions} places HLE action rates alongside their GPQA
counterparts.

\begin{table}[t]
\centering
\small
\caption{%
  Canonical action rates on GPQA-Diamond versus HLE.
  Each value = fraction of trajectories in the (harness, outcome) cell
  containing $\ge$1 event of that action type.
  GPQA values from event-log extraction; HLE values from keyword heuristics on
  trajectory text (DirectLLM) and normalized trace step types (OpenCode,
  ZeroClaw).
}
\label{tab:hle-cross-actions}
\setlength{\tabcolsep}{3pt}
\begin{tabular}{lllrrrrr}
\toprule
Bench. & Harness & Out. & Ans. & Reas. & Tool & Verify & Plan & Recov. \\
\midrule
\multirow{2}{*}{GPQA}
& OpenCode & C & 1.00 & 0.93 & \textbf{0.55} & 0.59 & 0.30 & 0.00 \\
& OpenCode & W & 1.00 & 0.92 & \textbf{0.15} & 0.49 & 0.32 & 0.00 \\
\midrule
\multirow{4}{*}{HLE}
& DirectLLM & C & 1.00 & 1.00 & 0.00 & 0.99 & 0.58 & \textbf{1.00} \\
& DirectLLM & W & 0.86 & 1.00 & 0.01 & 0.97 & 0.47 & \textbf{0.99} \\
\cmidrule{2-8}
& OpenCode & C & 1.00 & 1.00 & \textbf{0.23} & 1.00 & 0.56 & \textbf{0.70} \\
& OpenCode & W & 1.00 & 1.00 & \textbf{0.25} & 1.00 & 0.62 & \textbf{0.69} \\
\bottomrule
\end{tabular}
\end{table}

\findingbox{The tool-use rate gap between correct and wrong trajectories is
large on GPQA-Diamond (3.7$\times$ for OpenCode) and absent on HLE
(0.9$\times$ for OpenCode).
Tool activation is an outcome-discriminative signal only when the benchmark
contains executable subproblems.
Recovery and verification rates are uniformly higher on HLE across every
harness, reflecting the ARM-level finding that the model navigates genuine
knowledge uncertainty.}

On GPQA, OpenCode invokes tools in 55\% of correct versus 15\% of wrong
trajectories, a 3.7$\times$ outcome-conditioned gap.
On HLE, the same harness invokes tools in 23\% of correct and 25\% of wrong
trajectories: the gap collapses to 0.9$\times$.
HLE tool calls are exploratory (fact-checking, information retrieval) rather
than computational; they neither anchor the answer nor predict its correctness.

Two additional patterns emerge only in the cross-benchmark view.
First, recovery rates are uniformly higher on HLE: OpenCode rises from 0.00
(GPQA) to 0.70 (HLE); ZeroClaw from 0.00 to 0.71 to 0.73; DirectLLM HLE
reaches 0.99.
This reflects the ARM-level finding that RR is the dominant non-reason mode
on HLE.
Second, GPQA OpenClaw's plan explosion (0.16 correct to 0.31 wrong,
Table~\ref{tab:gpqa-action-rates}) does not appear on HLE, where plan rates
are uniformly moderate (0.47 to 0.73) and outcome-invariant.
The re-planning signature is specific to the combination of an open-loop
harness and a benchmark where PD is the natural cognitive default.


%%% ─────────────────────────────────────────────────────────────────────────
\subsection{AIME 2026}
\label{appendix:macro-aime}

We evaluate Qwen3.5-27B on the 30-problem AIME 2026 competition set under two
sampling regimes.
The \emph{k=1} regime runs a single sample per problem with an unrestricted
token budget (the full 131k-token context window); this serves as the
single-pass reference.
The \emph{k=32} regime runs 32 independent samples per problem under a
600-second per-sample wall-clock budget, applied to DirectLLM and all three
agent harnesses.
We report \emph{pass@1} (mean per-sample accuracy across all $30{\times}32$
samples), \emph{pass@32} (fraction of problems where at least one of 32
samples is correct), and \emph{no-prediction rate} (fraction of samples
that exhaust their budget before emitting any answer).
Table~\ref{tab:aime-overview} summarizes all five systems.

Two contrasts define the dataset.
DirectLLM k=1 achieves 83.3\% pass@1 through extended single-pass chains
averaging 24k tokens on correct problems and 95k on wrong ones, with 40\% of
wrong samples hitting the 131k-token cap.
Shifting to k=32 with a 600-second budget reduces DirectLLM's pass@1 to
30.6\% and generates 69\% silent timeouts; the same model, constrained,
produces no answer on more than two-thirds of its attempts.
Agent harnesses at k=32 achieve higher pass@1 (41--46\%) and substantially
higher pass@32 (73--80\%) than constrained DirectLLM, recovering through
parallel exploration what they lose in per-sample depth.

\begin{table}[t]
\centering
\small
\caption{%
  System-level overview on AIME 2026 ($n=30$, Qwen3.5-27B).
  DirectLLM (k=1) uses a single sample with full token budget.
  All k=32 systems use 32 samples per problem under a 600-second per-sample
  wall-clock budget.
  pass@1 = mean per-sample accuracy across all $30{\times}k$ samples.
  pass@32 = fraction of problems solved by at least one sample.
  No-pred\,\% = fraction of samples producing no extractable answer.
  Tok\,(C)/Tok\,(W) for DirectLLM k=1 from single-sample logprob data.
}
\label{tab:aime-overview}
\setlength{\tabcolsep}{3.5pt}
\begin{tabular}{lrrrrrr}
\toprule
System & pass@1\,\% & pass@32\,\% & No-pred\,\% & Tok (C) & Tok (W) & W/C \\
\midrule
DirectLLM (k=1)  & \textbf{83.3} & \textbf{96.7} &  0.0 & 24{,}304 &  94{,}950 & $3.9{\times}$ \\
\midrule
DirectLLM (k=32) & 30.6 & 53.3 & 69.4 & 10{,}214 & 12{,}623 & $1.2{\times}$ \\
OpenClaw  (k=32) & 40.9 & \textbf{80.0} & 58.1 &  2{,}613 &  7{,}661 & $2.9{\times}$ \\
OpenCode  (k=32) & \textbf{45.8} & 76.7 & 50.8 &  1{,}772 &  6{,}029 & $3.4{\times}$ \\
ZeroClaw  (k=32) & 41.1 & 73.3 & 50.3 &  2{,}955 &  3{,}333 & $1.1{\times}$ \\
\bottomrule
\end{tabular}
\end{table}

\findingbox{Extended per-sample thinking budget is the primary accuracy
driver for mathematical competition problems; constraining the budget
converts the dominant failure mode from a wrong answer to the absence of
any answer.}

DirectLLM k=1 wrong trajectories average 95k tokens and 40\% exhaust the
context window entirely: the model is genuinely attempting a solution and
running out of space.
This is not pathological length inflation; it reflects the difficulty of
competition mathematics, where extended derivations and multiple revisited
subproblems are necessary.
Under the 600-second k=32 budget, the same model achieves only 30.6\%
pass@1 and produces no answer in 69\% of samples, the highest no-prediction
rate of any system.
Table~\ref{tab:aime-budget} shows that budget compression does not produce
shorter wrong trajectories but produces far more empty outputs; the model
starts a derivation, runs out of time mid-chain, and emits nothing.

\textbf{Case: aime\_28 and aime\_24.}
Both problems are unsolved by any k=32 system across all 32 samples.
DirectLLM k=1 also fails, generating 131k-token attempts that exhaust the
full context window.
These represent the model knowledge ceiling: scaffolding and sampling cannot
compensate when the underlying derivation exceeds the model's capability.

\begin{table}[t]
\centering
\small
\caption{%
  DirectLLM budget scaling on AIME 2026.
  k=1 uses single-sample unrestricted budget; k=32 uses 600-second per-sample
  wall-clock limit.
  No-pred\,\% = fraction of samples producing no extractable answer.
  Cap\,\% = fraction of samples reaching the 131k-token context limit.
}
\label{tab:aime-budget}
\setlength{\tabcolsep}{4pt}
\begin{tabular}{lrrrrrr}
\toprule
System & pass@1\,\% & No-pred\,\% & Tok\,(C) & Tok\,(W) & W/C & Cap\,\%(W) \\
\midrule
DirectLLM (k=1)  & 83.3 &  0.0 & 24{,}304 & 94{,}950 & $3.9{\times}$ & 40 \\
DirectLLM (k=32) & 30.6 & 69.4 & 10{,}214 & 12{,}623 & $1.2{\times}$ &  0 \\
\bottomrule
\end{tabular}
\end{table}

\findingbox{Silent budget exhaustion is the dominant failure mode across all
systems at fixed compute; wrong answers are a minority of failures, making
pass@1 an incomplete diagnostic of per-sample competence.}

Among k=32 samples that fail, most produce no answer at all rather than a
wrong one.
DirectLLM k=32 has 69\% no-prediction; agent harnesses range from 50\% to
58\%.
This means that traditional accuracy metrics undercount failures: a sample
counted as ``wrong'' is a sample that committed and erred; a silent timeout
is counted identically but represents a different failure mode, one where the
model cannot complete the derivation within the allocated budget rather than
completing it incorrectly.
Table~\ref{tab:aime-nopred} shows that no-prediction rates are consistent
across outcome labels and do not correlate with problem-level difficulty
in the same way that wrong-answer rates do.

\begin{table}[t]
\centering
\small
\caption{%
  No-prediction rates for k=32 systems on AIME 2026.
  No-pred\,\% = fraction of all $30{\times}32$ samples producing no
  extractable answer before budget exhaustion.
  Tasks never solved = problems where all 32 samples produce no correct
  answer (pass@32 = 0).
}
\label{tab:aime-nopred}
\setlength{\tabcolsep}{4pt}
\begin{tabular}{lrrr}
\toprule
System & No-pred\,\% & Tasks never solved & pass@32\,\% \\
\midrule
DirectLLM (k=32) & 69.4 & 14 / 30 & 53.3 \\
OpenClaw  (k=32) & 58.1 &  6 / 30 & 80.0 \\
OpenCode  (k=32) & 50.8 &  7 / 30 & 76.7 \\
ZeroClaw  (k=32) & 50.3 &  8 / 30 & 73.3 \\
\bottomrule
\end{tabular}
\end{table}

\findingbox{With multiple independent samples at fixed total compute, agent
harnesses outperform unconstrained direct inference on the pass@k metric;
parallel exploration of different solution paths substitutes for the depth
that budget constraints remove from each individual sample.}

At equal total compute (32 samples at 600 seconds each), agent harnesses
achieve pass@32 of 73--80\% versus 53\% for DirectLLM.
The advantage comes from diversity: each sample in an agent harness may
invoke different tools, decompose the problem differently, or attempt a
different approach; if one of 32 approaches reaches the correct answer, the
task is solved.
DirectLLM k=32 samples are less diverse because they lack tool scaffolding
and tend to follow similar reasoning paths, producing a higher no-prediction
rate and a lower ceiling per additional sample.
The crossover between per-sample depth and multi-sample breadth occurs
somewhere between k=1 (DirectLLM leads by 37 percentage points on pass@1)
and k=32 (agent harnesses lead by 20 to 27 points on pass@32).

\textbf{Case: aime\_10 (combinatorics).}
DirectLLM k=1 generates 74k tokens and arrives at the wrong answer.
Under k=32, DirectLLM produces 30 empty samples and 2 wrong answers across
all 32 attempts.
OpenClaw k=32 solves the problem in one of its 32 samples through a different
decomposition, contributing to its 80\% pass@32.

\findingbox{The compute-accuracy scaling relationship plateaus at the model
knowledge ceiling; problems that no system solves across 32 samples
represent a hard boundary that additional scaffolding and sampling cannot
cross.}

Table~\ref{tab:aime-nopred} shows that DirectLLM k=32 leaves 14 of 30
problems unsolved across all 32 samples; the three agent harnesses leave
6 to 8 unsolved.
The three tasks that remain unsolved by every k=32 system (aime\_14,
aime\_15, aime\_28) are also among the hardest for DirectLLM k=1, with
131k-token attempts that exhaust the full context window.
These tasks represent the intersection of high problem difficulty and current
model capability limits: the underlying derivation exceeds what the model can
reliably complete, and no amount of harness scaffolding or repeated sampling
can compensate.
The gap between the oracle ceiling (any system solves the problem) and the
single-best-harness ceiling measures the recoverable portion of this limit;
on AIME, that gap shrinks as problem difficulty pushes more tasks past the
model knowledge boundary.
